\subsection{笛卡尔树}

构建时间复杂度为 $\mathcal{O}(n)$。

由下标区间 $[s,n)$ 的值建笛卡尔树：中序遍历恰为下标顺序，同时满足堆性质。堆方向由比较器 \texttt{CMP} 决定——默认 \texttt{std::less} 建出“大根”树（根为区间内值最大的下标，父值 $\ge$ 子值）；要建满足 $a_{\operatorname{parent}(u)}<a_u$ 的小根笛卡尔树则传入 \texttt{std::greater<T>}。

返回 $bst[u]=\{l,r\}$：下标 $u$ 结点的左/右儿子下标，无则 $-1$；左儿子下标总小于 $u$、右儿子大于 $u$（中序即下标序）。整棵树的根是唯一不为他人儿子的下标，即区间的最大（或最小）值所在位置。常见用途：区间 RMQ 转 LCA、笛卡尔树 DP、最大子矩形等。

\begin{lstlisting}
using p32 = std::array<int,2>;
template<typename T,typename CMP = std::less<T>>
std::vector<p32> Cartesian(std::vector<T>& arr,int s = 0,CMP compare = CMP()){
    std::vector<int> stk;
    std::vector<p32> bst(arr.size(),{-1,-1});
    for (int i = s;i < arr.size();i ++){
        while(stk.size() && compare(arr[stk.back()],arr[i])){
            bst[i][0] = stk.back();
            stk.pop_back();
        }
        if (stk.size()) bst[stk.back()][1] = i;
        stk.push_back(i);
    }
    return bst;
}
\end{lstlisting}